
\documentclass[12pt]{article}
\usepackage[utf8]{inputenc}
\usepackage{amsmath, amssymb, amsthm}
\usepackage{graphicx}
\usepackage{hyperref}
\usepackage{physics}
\usepackage{listings}
\usepackage{xcolor}
\usepackage{geometry}
\geometry{margin=1in}

\definecolor{codebg}{rgb}{0.95,0.95,0.95}
\lstset{
  backgroundcolor=\color{codebg},
  basicstyle=\ttfamily\footnotesize,
  breaklines=true,
  frame=single,
  postbreak=\mbox{\textcolor{red}{$\hookrightarrow$}\space}
}

\title{Lecture Notes on Quantum Boltzmann Machines}
\author{Prepared for Graduate Physics Students}
\date{\today}

\begin{document}

\maketitle

\tableofcontents
\newpage

\section{Introduction}

Quantum Boltzmann Machines (QBMs) are quantum generalizations of classical Boltzmann machines—stochastic neural networks used for modeling probability distributions. They aim to exploit quantum phenomena such as superposition and entanglement to potentially outperform their classical counterparts.

\section{Classical Boltzmann Machines}

\subsection{Energy-Based Models}

A classical Boltzmann machine is defined by an energy function:
\[
E(v, h) = -\sum_i b_i v_i - \sum_j c_j h_j - \sum_{i,j} v_i W_{ij} h_j
\]
where \( v \) are visible units, \( h \) are hidden units, \( W \) are weights, and \( b_i, c_j \) are biases.

\subsection{Probability Distribution}

The probability of a state is given by the Boltzmann distribution:
\[
P(v, h) = \frac{e^{-E(v,h)}}{Z}, \quad Z = \sum_{v,h} e^{-E(v,h)}
\]

\section{Quantum Boltzmann Machines}

\subsection{Quantum Generalization}

In the quantum case, the state of the system is described by a density matrix \( \rho \), and the energy is defined via a Hamiltonian \( H \):
\[
\rho = \frac{e^{-\beta H}}{Z}, \quad Z = \Tr(e^{-\beta H})
\]

\subsection{Hamiltonian Formulation}

We define a Hamiltonian over qubits:
\[
H = -\sum_i b_i Z_i - \sum_{i<j} W_{ij} Z_i Z_j + \sum_i \Delta_i X_i
\]
where \( Z_i, X_i \) are Pauli operators. The transverse field term \( \Delta_i X_i \) introduces quantum effects.

\section{Training Quantum Boltzmann Machines}

\subsection{Objective Function}

Training aims to minimize the KL divergence between the model distribution and data:
\[
\mathcal{L} = D_{KL}(P_{\text{data}} || P_{\text{model}})
\]

\subsection{Gradient Estimation}

The gradient of the loss involves expectation values:
\[
\partial_\theta \mathcal{L} = \expval{\partial_\theta H}_{\text{data}} - \expval{\partial_\theta H}_{\text{model}}
\]

\section{Implementation with PennyLane}

\subsection{Quantum Circuit Ansatz}

A variational circuit can be used to simulate the thermal state of the QBM. Here's a basic setup using PennyLane:

\begin{lstlisting}[language=Python]
import pennylane as qml
from pennylane import numpy as np

n_qubits = 4
dev = qml.device("default.qubit", wires=n_qubits)

@qml.qnode(dev)
def qbm_ansatz(params):
    for i in range(n_qubits):
        qml.RY(params[i], wires=i)
    for i in range(n_qubits - 1):
        qml.CNOT(wires=[i, i+1])
    return qml.expval(qml.PauliZ(0))
\end{lstlisting}

\subsection{Training Loop}

\begin{lstlisting}[language=Python]
from pennylane.optimize import GradientDescentOptimizer

params = np.random.randn(n_qubits, requires_grad=True)
opt = GradientDescentOptimizer(stepsize=0.1)

for i in range(100):
    params = opt.step(lambda p: -qbm_ansatz(p), params)
\end{lstlisting}

\section{Quantum Sampling}

Sampling from a quantum Boltzmann distribution is nontrivial. One approach is to use Quantum Annealing or Quantum Imaginary Time Evolution.

\section{Applications and Research Directions}

\begin{itemize}
  \item Quantum-enhanced generative modeling
  \item Quantum unsupervised learning
  \item Quantum advantage in statistical modeling
\end{itemize}

\section{Conclusion}

Quantum Boltzmann Machines represent a promising frontier in quantum machine learning, offering new computational paradigms through quantum statistics and variational principles.

\end{document}
